\begin{abstract}
Weighted Gerchberg--Saxton holograms equalize the traps of an optical tweezer array for the nominal optical system, and the aberrations that remain after a wavefront correction then change the trap intensities in ways the hologram was never asked to resist. We minimize the expected trap-intensity cost over a Gaussian prior on residual Zernike aberrations (Noll 4 to 11), replace the expectation by sixteen sigma points so that a conjugate-gradient descent from phase-fixed weighted Gerchberg--Saxton sees a noise-free cost with a closed-form gradient, and derive that the sigma points match the expectation through second order and that the robust cost adds the first-order intensity variance that the aberrations induce. On a $10\times10$ array in a scalar Fourier model, over an ensemble of \parNens{} draws of \parAb{} waves per mode, the 95th percentile of the relative standard deviation of the trap intensities is a median of \ratioMed{} times lower than for phase-fixed weighted Gerchberg--Saxton (\ratioMin{} to \ratioMax{} over three seeds; \ratioAllMed{} over \nSeedsAll{} seeds, interval \ratioAllCiLo{} to \ratioAllCiHi) at \effRatioMed{} of its efficiency. The same solver on the nominal cost gives \ablNominalMed, so the expected cost causes the gain; sixteen random draws give \ablRandSixteenMed. The gain falls to \misEMed{} at twice the design amplitude and \misGMed{} when modes outside the design are present, the first-order model overstates the gain, and for random displacements of the traps the acceptance configuration gives \ptAccMed{}, while a displacement prior optimized alone for 100 steps gives \poBLong{} at 0.10 focal bins at \poBLongEff{} of the efficiency. The results are simulations without pixel crosstalk and not hardware measurements.
\end{abstract}

\noindent\textbf{Keywords:} holographic optical tweezers, weighted Gerchberg--Saxton, residual aberrations, expected-cost optimization, sigma points, Zernike polynomials.
